% 与主文件同步的模板第2部分；使用主文件的宏定义。
\subsection{项目的研究内容、研究目标及拟解决的关键科学问题}

\subsubsection{研究内容}

本项目面向柔性桁架支撑下机器人的精确运动与模块组装，研究机器人与支撑结构的耦合动力学及受约束运动控制，开展动力学失配条件下的残差补偿与柔顺性能保持研究，进而解决目标位姿不确定条件下的接触对准与柔顺组装问题，并通过地面试验验证控制方法的有效性与适用范围。具体研究内容如下。

\topic{（1）柔性支撑机械臂的耦合动力学与受约束运动控制}

机器人依附柔性桁架作业时，机械臂运动产生的反作用载荷经基座传递至桁架，激发结构振动；桁架振动又通过基座运动作用于机械臂，改变末端位姿与动态响应，形成机器人运动与结构振动的双向耦合。本项目从机器人安装界面的运动与受力出发，研究基座、关节与末端之间的运动及载荷传递关系，揭示基座振动、作业载荷和机器人运动状态对末端响应的影响规律，将支撑作用表征为安装端运动与力和力矩，建立面向实时运动控制的耦合动力学描述。

面向位姿保持、轨迹跟踪及接触作业，研究关节运动范围与驱动力矩受限条件下的位置与姿态协同控制方法。分析基座振动和驱动约束对末端运动及柔顺响应的影响，研究运动精度、接触柔顺性与安装端反作用载荷之间的协调机制，形成振动扰动下的受约束精确运动控制方法，明确其在不同运动速度、作业载荷和基座振动条件下的适用范围。

\topic{（2）动力学失配下的运动精度补偿与柔顺性能保持}

针对惯性参数不确定、摩擦与负载变化及执行时延引起的动力学失配，研究模型与执行偏差对末端运动精度和安装端载荷的影响，分析不同偏差来源在机器人运动过程中的作用特征及变化规律。结合名义模型与实际运动响应，研究面向柔性支撑机器人的动力学残差表征与学习补偿方法，明确可由关节驱动有效补偿的偏差及其所需的控制作用。

围绕精度提升与柔顺性能保持，研究残差补偿与名义运动控制的协同关系，揭示补偿作用对末端阻抗响应及安装端反作用载荷的影响。重点分析驱动约束、基座振动和作业载荷变化对补偿效果的制约，统筹末端跟踪精度、接触柔顺性与补偿作用的平滑性，形成适应基座振动的残差补偿方法，并明确其在不同振动与载荷条件下的性能变化及适用范围。

\topic{（3）目标位姿不确定条件下的接触对准与柔顺组装}

针对目标位姿估计偏差与支撑运动共同作用下的接口对准问题，研究接口相对位姿、接触状态与接触载荷之间的关联，明确接触信息对位置和姿态偏差的辨识能力。分析支撑运动对接触几何判断及目标位姿估计的影响，研究基于接触信息的目标位姿修正方法，提高不确定条件下的接口对准能力。

面向接口对准、插接及到位保持过程，研究位姿修正与柔顺运动的协调方法，实现横向及姿态纠偏与轴向接近的协同控制。考察目标估计偏差、运动执行误差和局部接触约束对接触载荷及安装端受力的影响，揭示对准停滞、载荷持续增大等现象的形成条件。在典型接口和既定连接构型下，研究基座振动条件下兼顾对准精度与接触载荷的稳定组装条件，确定可完成对接的初始位姿偏差范围。

\topic{（4）柔性支撑机器人运动与组装的地面试验验证}

围绕耦合运动控制、残差补偿与柔顺组装的验证需求，采用KUKA末端悬挂Franka的实验构型，由KUKA施加基座运动激励、Franka执行运动与组装任务，模拟在轨柔性支撑振动引起的安装端运动。建立基座与作业端运动、安装端载荷及接触载荷的同步测量方法，分析地面重力、悬挂连接与激励执行偏差对试验结果的影响，明确地面基座激励与在轨支撑运动的对应范围，制定可重复的试验流程。

开展位姿保持、轨迹跟踪和典型接口组装的分项与集成试验，通过不同基座振动、作业载荷及初始位姿偏差条件下的对照和重复验证，分别评价受约束运动控制、残差补偿与接触纠偏的效果。在可比的基座激励下，综合分析末端运动误差、安装端反作用载荷、接触载荷和组装到位误差，检验各控制环节协同作用下的作业性能，完成柔性支撑振动环境下机器人运动与组装的地面原理验证。

\subsubsection{研究目标}

总体目标是阐明柔性支撑、动力学失配与接触状态变化对机器人运动和柔顺响应的影响，建立面向模块组装的精确运动与柔顺控制方法，并通过地面试验验证多源不确定性条件下的作业能力。与四项研究内容对应，具体目标如下。

\topic{（1）建立柔性支撑条件下的受约束运动控制方法}

明确安装界面运动与作业端响应之间的动力学联系，形成兼顾位置、姿态及驱动约束的运动控制方法，实现代表性基座振动条件下的位姿保持与轨迹跟踪，并给出运动精度、柔顺性和安装端反作用载荷之间的协调依据。

\topic{（2）提升动力学失配条件下的运动精度与柔顺性能}

减轻惯性、摩擦、负载与执行时序偏差对控制性能的影响，明确偏差补偿对末端精度和安装端载荷的作用范围。在相同基座激励、任务参考与驱动约束下，以较未引入补偿的名义控制的位置跟踪均方根误差降低不小于20\%、振动激励下位姿保持阶段的末端位置均方根误差降低不小于25\%为拟定目标，同时评价姿态误差、柔顺响应、安装端载荷和力矩变化，验证补偿效果及其代价。

\topic{（3）形成目标位姿不确定条件下的柔顺组装能力}

明确接触信息对目标位置和姿态偏差的约束能力，建立基座振动条件下兼顾目标对准与接触载荷的柔顺组装方法。在预定初始位姿偏差范围内实现典型接口的对准、插接及到位保持，明确目标偏差修正与柔顺执行的配合关系及稳定作业条件。

\topic{（4）完成运动控制与典型接口组装的地面原理验证}

完善KUKA末端悬挂Franka的基座振动模拟平台，配置典型接口与同步测量系统，完成3组基座振动工况与2种作业载荷组合下的运动对照，以及不同初始位姿偏差下的组装验证。形成可重复的试验流程与多工况数据，分别评价控制、补偿和接触纠偏的贡献，报告集成系统的运动精度、安装端载荷、组装到位精度及接触载荷，明确方法的适用范围。

改善比例为拟定目标，采用预定工况和评价窗口，同时报告绝对误差与重复试验离散程度。对接以独立测量的真实目标坐标和接口容差评价，分别报告到位精度、接触载荷及可完成对接的偏差范围。

\subsubsection{拟解决的关键科学问题}

\topic{（1）柔性支撑与模型失配共同作用下末端阻抗响应的形成机制}

柔性支撑运动和机械臂内部模型失配都可表现为末端偏差，但前者来自机器人与结构之间的真实动力学相互作用，后者体现为名义模型与实际执行之间的偏离。二者对驱动力矩的需求并不相同，若不加区分地学习和抵消末端误差，可能在改善局部跟踪的同时改变支撑受力和系统柔顺性。另一方面，关节与力矩约束使期望阻抗响应不能始终被准确实现，残差补偿的实际效果随约束状态而变化。

因此，需要揭示支撑作用、内部失配及受约束力矩分配共同影响末端响应的机制，确定驱动能够有效补偿的偏差及补偿对安装端反作用载荷的影响。其核心在于明确学习补偿改善振动扰动下运动精度的条件，以及保持所需柔顺性、限制附加反作用载荷的控制边界，为柔性结构上的学习增强控制提供依据。

\topic{（2）柔性支撑条件下接触几何修正与阻抗响应的协调机制}

对接时的接触位姿和载荷同时受目标偏差、支撑运动与机器人柔顺响应影响，较小的接触力不一定对应正确的装配位置，局部接触数据也未必能充分约束所有位姿方向。利用接触观测更新目标或运动参考，又会改变接触状态、关节作用和安装端载荷，使几何估计与运动控制相互影响。

需要明确接触信息约束目标偏差的能力及可信程度，揭示基座振动条件下参考变化经实际阻抗控制传递至接触载荷和安装端受力的规律，使参考调整既能利用接触反馈纠偏，又保持对目标位姿的约束，从而协调振动扰动下的对准精度与柔顺性，实现典型工况下的稳定对接。

